%DS0



\section{ĐỀ SỐ \thesection }
\setcounter{bt}{0}
%%Câu 1.
\begin{bt} Thực hiện phép tính:
	\begin{listEX}[3]
		\item $35-50+(-15)$; 
		\item $2^{3} \cdot 3^{2}-35: 5+2023^{0}$;
		\item $98+2 \cdot\left[\left(3^{2}+5\right): 7+1\right]$.
	\end{listEX}
	
\loigiai{a) $-30$ \hfill b) 66 \hfill c) 104}

\end{bt}


%%Câu 2. 
\begin{bt} Tìm $x$ biết:
	\begin{listEX}[3]
		\item $x+14=11+2x$;
		\item $3\cdot x-9=-36$;
		\item $(7-x) - (25+7 ) = -25$.
	\end{listEX}
	
	\loigiai{a) $3$ \hfill b) $-9$ \hfill c) 0}
	
	\end{bt}

%%Câu 3.
\begin{bt} Lớp 6A có $24$ học sinh nam và $20$ học sinh nữ. Trong buổi thi nghi thức đội, giáo viên chủ nhiệm lớp 6A đã yêu cầu lớp trưởng chia các bạn thành các nhóm sao cho giữa các nhóm phải có số học sinh nam bằng nhau, số học sinh nữ bằng nhau. Hỏi rằng lớp trưởng chia được nhiều nhất bao nhiêu nhóm?
	
	
\loigiai{$4$ nhóm}
\end{bt}


%%Câu 4.
\begin{bt} Bạn Yến được mẹ cho $500\,000$ đồng để đi nhà sách. Bạn đã mua $12$ quyển vở giá $10\,500$ đồng/quyển, $1$ bộ sách giáo khoa giá $190\,000$ đồng/bộ, $1$ bút chì giá $36\,000$ đồng/cây và $6$ bút gel với giá $8\,000$ đồng/cây.
	\begin{listEX}
		\item Hỏi bạn Yến đã mua hết bao nhiêu tiền?
		\item Với số tiền còn lại, bạn có thể mua được mấy ly kem Mixue biết giá 1 ly kem là $25\,000$ đồng/ly?
	\end{listEX}
	
\loigiai{a) $400\,000$ đồng \hfill b) $4$ ly kem}

	\end{bt}
	

%%Câu 5.
\begin{bt} Khu vườn hình chữ nhật có chiều rộng $12$ m và chiều dài gấp đôi chiều rộng. Tính diện tích của khu vườn.
	
\loigiai{288 m$^2$}
\end{bt}


%%Câu 6.
\begin{bt} Biểu đồ tranh dưới đây thể hiện số cây thông trang trí bán được trong 4 tuần của tháng 12 tại một cửa hàng.
	\begin{center}
		\def\icon{\faTree}
		\def\hicon{\tikz{\clip(-0.18,-0.17)rectangle(0,0.4);\path(0,0)node{\faTree}}}
		\def\vetranh#1{\foreach \dem in{1,...,#1}{\icon \ }}{
			\begin{tabular}{|c|p{6cm}<{}|}\hline 
				\multicolumn{2}{|c|}{\bf Số cây thông trang trí trong tháng 12}\\ \hline 
				\bf Tuần &\multicolumn{1}{c|}{\bf Số cây thông}\\ \hline 
				Tuần 1 &\vetranh{3}\hicon\\ \hline 
				Tuần 2 &\vetranh{5} \\ \hline 
				Tuần 3 &\vetranh{6}\hicon \\ \hline 
				Tuần 4 &\vetranh{8} \\ \hline 
			\end{tabular}\\
			(\icon = 10 xe cây thông;        \hicon = 5 cây thông)}
	\end{center}
		\begin{listEX}
			\item Lập bảng thống kê biểu diễn số cây thông trang trí bán được trong 4 tuần của tháng 12.
			\item Tính tổng số cây thông trang trí bán được trong 4 tuần của tháng 12.
		\end{listEX}
	
\loigiai{a) Học sinh tự lặp bảng \hfill b) 230 cây thông}

\end{bt}

%%Câu 7.
\begin{bt} Cho biểu thức: $A=7+7^{2}+7^{3}+7^{4}+7^{5}+\cdots+7^{2022}+7^{2023}$\\
Tìm $x$ để $6A+7=7^{x}$
\loigiai{$x=2024$}
\end{bt}
